\section{A worked example: exact GBC for the normal model}

Let $\theta\sim N(0,\sigma_0^2)$ and $y\mid\theta\sim N(\theta,\sigma^2)$. The posterior is $N(by,b\sigma^2)$ with $b = \sigma_0^2/(\sigma_0^2+\sigma^2)$, and its quantile function is
\[
H(y,u) = by+\sqrt{b\sigma^2}\,\Phi^{-1}(u),
\]
where $\Phi$ is the standard normal distribution function. The map is affine in $y$ for each quantile level. A quantile regression of simulated $\theta_i$ on simulated $y_i$ with the check loss at level $u$, using linear functions of $y$, recovers $H$ exactly in the limit of many simulations; the error at sample size $N$ is of order $N^{-1/2}$ by standard quantile regression asymptotics. The efficiency of learning depends on how well the function class captures the conditional quantile function.

\section{Sinkhorn's algorithm and the Hilbert metric}
\label{sec:sinkhorn_hilbert}

Birkhoff's theory of positive operators gives a proof that Sinkhorn's iteration converges geometrically. For positive vectors define Hilbert's projective metric~\cite{hilbert1895}
\[
d_H(a,a') = \log\max_i\frac{a_i}{a'_i}+\log\max_j\frac{a'_j}{a_j},
\]
a metric on rays of the positive cone. Birkhoff showed in 1957~\cite{birkhoff1957} that a matrix $G$ with positive entries is a strict contraction for $d_H$, with contraction ratio $\tanh(\Delta/4)$, where $\Delta$ is the projective diameter of the image of the cone. Each half step of Sinkhorn's algorithm applies $G$ or $G^\top$ followed by a coordinatewise division, which is an isometry for $d_H$, so the scaling vectors converge geometrically, with ratio $\tanh^2(\Delta/4)$ per iteration. Franklin and Lorenz used this argument in 1989~\cite{franklinlorenz1989}. For the Gaussian kernel $G_{ij} = e^{-|x_i-y_j|^2/2\varepsilon}$ one finds $\log(G_{ik}G_{jl}/G_{il}G_{jk}) = (x_i-x_j)(y_k-y_l)/\varepsilon$ for rows $i,j$ and columns $k,l$, so that $\Delta = L_xL_y/\varepsilon$ for grids of diameters $L_x$ and $L_y$, and the ratio $\tanh(\Delta/4) = 1-O(e^{-\Delta/2})$ approaches one as $\varepsilon\to0$. The ratio is a worst case bound, and it is often far from sharp when the unregularized problem has a Monge map~\cite[Section~4.2]{peyrecuturi2019}, as it does in Figure~\ref{fig:sinkhorn}. There the grids span $[-4,4]$ and the ratio is within $10^{-13}$ of one in all three panels, yet the $\ell^1$ error of the marginals falls below $10^{-12}$ after about thirty, two hundred and two thousand iterations. The count grows in proportion to $1/\varepsilon$, in agreement with the local rate $1-\kappa$, $\kappa\sim\varepsilon$, that is observed numerically when a smooth Monge map exists~\cite[Section~4.2]{peyrecuturi2019}. Small regularization is therefore slow. Altschuler, Weed and Rigollet~\cite[Theorem~2]{altschuler2017} proved a polynomial bound: for a positive matrix with sum $s$ and smallest entry $m$, the iteration reaches an $\ell^1$ marginal error $\delta$ after $O(\delta^{-2}\log(s/m))$ steps, and $\log(s/m)$ grows like $1/\varepsilon$ for the Gaussian kernel.

\section{The F\"ollmer process and stochastic interpolants}

A simple bridge between a Dirac mass and a target law $\mu$ on $\R^d$ is the F\"ollmer process~\cite{follmer1985,lehec2013}: Brownian motion on $[0,1]$ started at zero and conditioned to have law $\mu$ at time one. It is the Schr\"odinger bridge from $\delta_0$ to $\mu$, its law is $\int B^{0y}\mu(\dd y)$ by Theorem~\ref{thm:static_reduction} with $\varepsilon=1$, where $B^{0y}$ is the Brownian bridge from zero to $y$, and its drift is
\[
b(t,x) = \nabla\log\E\Big[\frac{\dd\mu}{\dd\gamma}\big(x+\sqrt{1-t}\,Z\big)\Big],
\]
where $\gamma$ is the standard Gaussian law and $Z$ standard Gaussian. The process has the smallest relative entropy with respect to Wiener measure among all processes from zero to $\mu$, namely $\KL(\mu\|\gamma)$. Sampling by simulating the F\"ollmer process requires the drift, which can be estimated by Monte Carlo or learned. Stochastic interpolants~\cite{albergo2023} and flow matching~\cite{lipman2023} generalize the idea: one prescribes an interpolation $X_t = \alpha_tX_0+\beta_tX_1$ between noise and data, learns the conditional expectation of its velocity given its position by regression, and integrates the resulting ordinary or stochastic differential equation. The regression target is a conditional expectation, so the learning problem is again of the GBC type: simulate pairs, regress, and use the regression to transport noise to the target.

\section{Approximate Bayesian computation and its limits}

Approximate Bayesian computation, introduced by Tavar\'e, Balding, Griffiths and Donnelly~\cite{tavare1997} and by Pritchard and coauthors~\cite{pritchard1999} in the 1990s, accepts simulated parameters whose simulated summary statistics fall within a tolerance of the observed ones. In the limit of zero tolerance and sufficient summaries it is exact, but its accuracy degrades with the dimension $d$ of the summary: with $n$ simulations the minimal mean squared error of the estimated posterior density is of order $n^{-4/(d+5)}$, for rejection as for regression adjustment~\cite[Remark~1]{blum2010}. For a given accuracy the number of simulations then grows exponentially in $d$, up to constants that also depend on $d$. GBC removes the tolerance but not the dependence on dimension. The minimax rate for estimating a $p$ times differentiable regression function of $d$ variables is $n^{-p/(2p+d)}$ in $L^2$~\cite{stone1982}, and the conditional quantile function is no easier to learn, since in a location model every conditional quantile differs from the regression function by a constant. GBC therefore relies on structure in the quantile map and on a low dimensional summary statistic of the data. Regression adjustment, introduced by Beaumont, Zhang and Balding~\cite{beaumont2002} in 2002, corrects accepted parameters by a regression on the summaries and is a first order version of the quantile map. The progression from rejection to regression adjustment to learned generators parallels the progression in kinetic theory from direct simulation of particles to moment closures to learned surrogate models: each step replaces sampling by estimation of a conditional expectation.

\section*{Exercises}
\begin{enumerate}
\item Show that the minimizer of the check loss is unique when the distribution function is strictly increasing.
\item Verify the identity $(I-KA)P(I-KA)^\top+KRK^\top = (I-KA)P$.
\item Let $\mu$ have no atoms and let $T = F_\nu^{-1}\circ F_\mu$. (a) Show that $T_\#\mu = \nu$, and that $T$ is optimal for the quadratic cost when $\mu$ is absolutely continuous and both measures have finite second moments, using Theorem~\ref{thm:brenier}. (b) Let $c$ be convex. Show that $c(x-y)+c(x'-y')\leq c(x-y')+c(x'-y)$ for $x<x'$ and $y<y'$, and deduce that for $x_1<\dots<x_n$ and $y_1<\dots<y_n$ the matching $x_i\leftrightarrow y_i$ minimizes $\sum_ic(x_i-y_{\sigma(i)})$ over the permutations $\sigma$. (c) For $c(z) = |z|$, $\mu$ uniform on $[0,1]$ and $\nu$ uniform on $[1,2]$, show that every coupling has the same cost, so that $T$ is optimal but not unique.
\item Implement Sinkhorn's algorithm for two discrete measures on a grid and check that the marginal errors decrease geometrically. Estimate the rate per iteration for several values of $\varepsilon$ and compare it with $\tanh^2(\Delta/4)$, where $\Delta = L_xL_y/\varepsilon$ for the kernel $e^{-|x-y|^2/2\varepsilon}$ and grids of diameters $L_x$ and $L_y$.
\item Use Theorem~\ref{thm:tweedie} to compute the posterior mean when the prior is uniform on $\{-1,1\}$ and the noise is $N(0,\sigma^2)$.
\end{enumerate}
